\subsection{The distinguished cover in the almost-ruled case}

The remaining almost-ruled alternative determines the particular
cover defined by the torsion line bundle, rather than only some
purely inseparable cover of the surface.

\begin{proposition}\label{prop:distinguished-ruled-cover}
Let $X$ be a smooth projective integral surface over an
algebraically closed field $k$ of characteristic $p>0$.
Let $Y\subset X$ be a smooth integral ample divisor, and let
$M\ne\OO_X$ have order $p$ with $M|_Y\simeq\OO_Y$.
Let $N$ be the normalization of the $\mu_p$-torsor defined by
a trivialization of $M^{\otimes p}$.
If a ruled surface $R$ admits a dominant purely inseparable
rational map of degree $p$ to $X$, then
\[
 k(R)=k(N)\quad\text{as extensions of }k(X).
\]
In particular $N$ is birationally ruled, and every smooth
projective model of its function field has Kodaira dimension
$-\infty$. A resolution of the rational map from a smooth
projective ruled model gives a proper birational morphism
to $N$ over $X$. The surface $X$ itself is not birationally ruled.
Moreover, if $k(N)=A(t)$ is a ruling presentation, then
\[
 k(X)\cap A=A^p,\qquad k(X)A=k(N).
\]
\end{proposition}

\begin{proof}
Choose local frames $e_i$ of $M$ and write $u_i\in\OO_X^*$
for their $p$th powers under the chosen trivialization.
On an overlap, $u_i=g_{ij}^p u_j$, so the forms $du_i/u_i$
glue to a regular global form $\omega$ on $X$.
It is nonzero. Indeed, for a rational frame $e$, let
$u\in K^*$ be its $p$th power, where $K=k(X)$.
If $du=0$, the perfect-field differential criterion
\cite[Tag~031W]{Stacks} gives $u=v^p$ with $v\in K$.
The rational section $e/v$ has $p$th power one.
Its divisor has every coefficient zero, so normality makes
it a nowhere-vanishing global section of $M$, a contradiction.
Thus $u\notin K^p$, and the torsor has the degree-$p$ field
\[
 E=K(u^{1/p}).
\]
Its locally free algebra injects into that generic field,
so the torsor is integral. Its normalization equals the
normalization of $X$ in $E$, since integrality is transitive.
This normalization is finite: apply the Nagata finiteness
criterion to $\Spec E\to X$
\cite[Lemma~29.54.14, Tag~0BAK]{Stacks}.

The image of $\omega|_Y$ in $\Omega^1_{Y/k}$ is zero.
A trivialization of $M|_Y$ has $p$th power differing from
the chosen trivialization by a global unit on the proper
integral curve $Y$, hence by an element of $k^*$.
Its logarithmic differential is zero.
The conormal sequence is
\[
 0\longrightarrow\OO_Y(-Y)\longrightarrow
 \Omega^1_{X/k}|_Y\longrightarrow\Omega^1_{Y/k}
 \longrightarrow0.
\]
Since $Y^2>0$, the left-hand line bundle has no nonzero
global section. Hence $\omega|_Y$ vanishes in the full
restricted cotangent bundle. There is consequently a
nonzero section
\[
 \theta\in H^0(X,\Omega^1_{X/k}(-Y)),
 \qquad s_Y\theta=\omega,
\]
where $s_Y$ is the canonical section of $\OO_X(Y)$.

Write the ruled function field as $k(C)(t)$, with $C$ a
smooth projective integral curve, and replace $R$ by
$R_0=C\times\PP^1$. Resolve $R_0\dashrightarrow X$ by
point blowups, obtaining $S\to R_0$ and $g:S\to X$
\cite[Tag~0C5H]{Stacks}. Point blowups of a regular surface
remain regular \cite[Lemma~54.3.2, Tag~0AGP]{Stacks}; since
$k$ is perfect, $S$ is smooth. It is projective, and
$[k(S):K]=p$.

Only finitely many ruling fibres contain images of the
blowup centres. A general fibre $F$ of $S\to C$ is therefore
an unmodified projective line, with
\[
 \Omega^1_{S/k}|_F\simeq\OO_{\PP^1}\oplus\OO_{\PP^1}(-2).
\]
The restriction $g|_F$ is nonconstant: a general ruling
fibre meets the dense quasi-finite locus of $g$ and cannot
be contracted. Ampleness of $Y$ gives
$e_F=\deg(g^*\OO_X(Y)|_F)>0$, without any separability
assumption on this curve map.

Pullback by the morphism $g$ gives a regular section
\[
 \eta\in H^0(S,\Omega^1_{S/k}(-g^*Y)).
\]
It restricts on each general ruling fibre to a section of
\[
 \OO_{\PP^1}(-e_F)\oplus\OO_{\PP^1}(-e_F-2),
\]
and is therefore zero there. These fibres fill a dense open
of the integral surface $S$, so $\eta=0$ and $g^*\omega=0$.
This uses positivity only on the general fibre; $g^*Y$
need not be ample on $S$.

At the generic point, $du/u=0$ in $\Omega^1_{k(S)/k}$.
The same perfect-field criterion gives a $p$th root of $u$
in $k(S)$, and hence a $K$-embedding $E\hookrightarrow k(S)$.
Both extensions have degree $p$, so they are equal.
The $p$th root is unique, and changing the frame or the
trivialization changes $u$ only by a $p$th power; thus this
is the distinguished extension defined by $M$.

The resulting birational map $S\dashrightarrow N$ over $X$
extends to a morphism. For an affine open $U=\Spec B$ of
$X$, let $B'$ be the integral closure of $B$ in $E$.
Every $b'\in B'$ is integral over every local ring
$\OO_{S,s}$ for $s\in g^{-1}(U)$, and belongs to it by
normality. These regular functions give
$g^{-1}(U)\to\Spec B'$ and glue over $X$.
The resulting map is birational by the field equality and
proper by its closed graph and properness of $g$.
Thus $N$ has a ruled function field.

For any smooth projective model $S'$ of that field, resolve
the birational map $R_0\dashrightarrow S'$ in the same way.
The resulting ruled blowup has no positive pluricanonical
sections: its canonical bundle has degree $-2$ on every
general unmodified ruling fibre. Pullback by the birational
morphism injects the pluricanonical sections of $S'$ into
those of the blowup. Hence every positive plurigenus of
$S'$ is zero.

If $X$ itself were birationally ruled, the argument with
$\theta$ could instead be applied to a birational map
$R_0\dashrightarrow X$. It would give $g^*\omega=0$, although
the map of function fields and hence of their differentials
is an isomorphism. This contradicts $\omega\ne0$.

Finally write $L=k(N)=A(t)$. The purely inseparable
degree-$p$ extension $L/K$ has exponent one, so
$L^p\subset K$ and $A^p\subset A\cap K\subset A$.
Since $A$ is the function field of a curve over the perfect
field $k$, one has $[A:A^p]=p$. For example, a separating
generator $x$ gives $[A:k(x)]=[A^p:k(x^p)]$, and the tower
through $k(x^p)\subset k(x)$ proves the assertion.
If $A\subset K$, then
$A(t^p)\subset K\subset A(t)$ and the degree-$p$ equality
forces $K=A(t^p)$, contradicting non-ruledness of $X$.
Thus $A\cap K=A^p$. The strict intermediate compositum
$KA$ in the prime-degree extension $L/K$ is all of $L$.
\end{proof}

The field intersection does not assert that the rational ruling
map on $X$ extends everywhere, or that its fibres are smooth.
The normalization may be singular away from the inverse image
of $Y$; resolving those singularities changes an ample pullback
into a divisor having degree zero on exceptional curves.
